%%%%%%%%%%%%%%%%%%%%%%%%%%%%%%%%%%%%%%%%%%%%%%%%%%%%%%%%%%%%%%%%%%%%%%%%%%%
\chapter{Retrieval Architectures}
\label{ch:retrieval-arch}
%%%%%%%%%%%%%%%%%%%%%%%%%%%%%%%%%%%%%%%%%%%%%%%%%%%%%%%%%%%%%%%%%%%%%%%%%%%

Embedding models can produce different kinds of representations. The architecture determines
what kind of index you build and how query-document similarity is computed.

%%%%%%%%%%%%%%%%%%%%%%%%%%%%%%%%%%%%%%%%%%%%%%%%%%%%%%%%%%%%%%%%%%%%%%%%%%%
\section{Dense Retrieval}
%%%%%%%%%%%%%%%%%%%%%%%%%%%%%%%%%%%%%%%%%%%%%%%%%%%%%%%%%%%%%%%%%%%%%%%%%%%

The standard approach. One vector per text.

\begin{Verbatim}
query  → [single vector]
doc    → [single vector]
score  = dot_product(query_vec, doc_vec)
\end{Verbatim}

Fast at query time. Index is compact. Works well when queries and documents are semantically
similar in the embedding space.

\textbf{Weakness:} a single vector must compress all meaning. Queries that need to match
multiple distinct aspects of a document can lose signal.

%%%%%%%%%%%%%%%%%%%%%%%%%%%%%%%%%%%%%%%%%%%%%%%%%%%%%%%%%%%%%%%%%%%%%%%%%%%
\section{Sparse Retrieval (SPLADE)}
%%%%%%%%%%%%%%%%%%%%%%%%%%%%%%%%%%%%%%%%%%%%%%%%%%%%%%%%%%%%%%%%%%%%%%%%%%%

Instead of a dense vector, the model outputs a sparse weighted term vocabulary---like a
learned BM25.

\begin{Verbatim}
query → {term: weight, term: weight, ...}  (most weights ≈0)
doc   → {term: weight, term: weight, ...}
score = sum of overlapping weighted terms
\end{Verbatim}

Good at exact and near-exact term matching. Complements dense retrieval on queries where
lexical specificity matters: version numbers, proper nouns, code tokens.

BGE-M3 includes a sparse retrieval head alongside its dense head, making hybrid pipelines
cheaper to operate.

%%%%%%%%%%%%%%%%%%%%%%%%%%%%%%%%%%%%%%%%%%%%%%%%%%%%%%%%%%%%%%%%%%%%%%%%%%%
\section{ColBERT --- Multi-vector / Late Interaction}
%%%%%%%%%%%%%%%%%%%%%%%%%%%%%%%%%%%%%%%%%%%%%%%%%%%%%%%%%%%%%%%%%%%%%%%%%%%

ColBERT keeps one vector per token rather than collapsing to a single vector.

\begin{Verbatim}
query: [what]  [is]  [the]  [capital]
          ↓      ↓     ↓       ↓
          q1     q2    q3      q4

doc:  [Paris] [is] [a] [city] [in] [France]
          ↓      ↓   ↓    ↓      ↓     ↓
          d1     d2  d3   d4     d5    d6
\end{Verbatim}

Similarity (MaxSim):

\begin{Verbatim}
score = sum over each qi of: max(dot(qi, dj) for all dj)
\end{Verbatim}

Each query token independently finds its best matching document token, then scores are summed.

\textbf{Why it matters:}

\begin{itemize}
    \item ``bank of a river'' and ``bank loan'' map to the same dense vector in ambiguous
          models; ColBERT keeps them distinct because \textit{river} and \textit{loan} match
          different document tokens
    \item Multi-aspect queries (``fast AND cheap AND local'') get each aspect separately
          matched rather than compressed into one average direction
    \item Exact token sensitivity: ``Python~3.11 migration'' won't match a document about
          ``Python~2.7'' because \textit{3.11} and \textit{2.7} pull in different directions
\end{itemize}

\textbf{Tradeoffs:}

\begin{center}
\begin{tabular}{llll}
\toprule
& \textbf{Dense} & \textbf{Sparse} & \textbf{ColBERT} \\
\midrule
Index size & small & medium & large (tokens $\times$ docs) \\
Query speed & very fast & fast & slower \\
Recall on complex queries & good & good for lexical & better \\
Typical use & first-stage retrieval & hybrid retrieval & reranking or dedicated index \\
\bottomrule
\end{tabular}
\end{center}

%%%%%%%%%%%%%%%%%%%%%%%%%%%%%%%%%%%%%%%%%%%%%%%%%%%%%%%%%%%%%%%%%%%%%%%%%%%
\section{Hybrid Retrieval}
%%%%%%%%%%%%%%%%%%%%%%%%%%%%%%%%%%%%%%%%%%%%%%%%%%%%%%%%%%%%%%%%%%%%%%%%%%%

Most production systems combine modes rather than choosing one:

\begin{enumerate}
    \item Dense retrieval $\to$ top-k candidates (fast, high recall)
    \item Sparse/ColBERT reranking $\to$ reorder top-k (slower, higher precision)
\end{enumerate}

BGE-M3 is notable because it runs all three modes (dense, sparse, ColBERT) from a single set
of weights, making hybrid pipelines cheaper to operate than maintaining separate models.

The typical local deployment for this repo uses dense retrieval for first-stage candidate
generation and reserves sparse or ColBERT for a second-stage reranking pass.

%%%%%%%%%%%%%%%%%%%%%%%%%%%%%%%%%%%%%%%%%%%%%%%%%%%%%%%%%%%%%%%%%%%%%%%%%%%
\section{Retrieval vs Re-ranking: The Two-Stage Pattern}
%%%%%%%%%%%%%%%%%%%%%%%%%%%%%%%%%%%%%%%%%%%%%%%%%%%%%%%%%%%%%%%%%%%%%%%%%%%

This distinction is one of the most important things to understand when choosing an embedding
stack.

\subsection{Stage 1: Ranking (Retrieval)}

Ranking is the initial ordering produced by the retriever. In an embedding pipeline:

\begin{enumerate}
    \item Embed the query
    \item Embed all documents or chunks ahead of time
    \item Compute similarity
    \item Return the top $k$ nearest items
\end{enumerate}

This stage is fast, scalable, and approximate. Document embeddings are precomputed; nearest-
neighbor search can be indexed efficiently.

\subsection{Stage 2: Re-ranking}

Re-ranking is a second pass over the candidate set returned by the retriever:

\begin{enumerate}
    \item Retriever returns top 20, 50, or 100 candidates
    \item Reranker reads the query plus each candidate together
    \item Reranker assigns a more exact relevance score
    \item Candidates are reordered before being passed to generation
\end{enumerate}

This stage is slower and more expensive, but usually more accurate, because the reranker can
reason over the direct query-document interaction rather than coarse vector proximity.

\subsection{Why Not Rerank Everything?}

Reranking does not scale to a large corpus. If you have 1M chunks, you can vector-search them
efficiently, but you do not want to run a cross-encoder reranker over all 1M. So the pattern
is: retrieve broadly, rerank narrowly.

\subsection{What Each Stage Optimizes For}

\begin{center}
\begin{tabular}{ll}
\toprule
\textbf{Stage} & \textbf{Goal} \\
\midrule
Embedding model & Maximize semantic recall at scale; good-enough ranking \\
Reranker & Maximize precision at the top of the list \\
\bottomrule
\end{tabular}
\end{center}

Embedding and reranking scores are not interchangeable. A strong embedding model that gets
the right document in the top 20 at high speed is more valuable than a slow but precise model
that cannot serve queries under load.

\subsection{Practical Retrieval Stack}

A strong modern local retrieval stack:

\begin{enumerate}
    \item Embedding model generates query and document vectors
    \item ANN index (HNSW) returns top 20--100 candidates (see Chapter~\ref{ch:ann})
    \item Reranker rescoring narrows to best top 5--10
    \item Only the highest-ranked chunks go into the final LLM prompt
\end{enumerate}

%%%%%%%%%%%%%%%%%%%%%%%%%%%%%%%%%%%%%%%%%%%%%%%%%%%%%%%%%%%%%%%%%%%%%%%%%%%
\section{Architecture Selection Guide}
%%%%%%%%%%%%%%%%%%%%%%%%%%%%%%%%%%%%%%%%%%%%%%%%%%%%%%%%%%%%%%%%%%%%%%%%%%%

\begin{description}
    \item[First local deployment] Use dense retrieval with Qwen3-Embedding-4B. Serve with
          \texttt{llama-server~--embedding~--pooling~last}. Add reranking if top-of-list
          precision becomes a bottleneck.
    \item[Hybrid search needed] Add BM25 alongside vector search; fuse results with RRF (see
          Chapter~\ref{ch:ranking}). memvid includes this by default in its \texttt{ask()} method.
    \item[ColBERT / multi-vector] Introduce BGE-M3 if you need multi-aspect query matching
          and are willing to manage a larger index. Not the first step.
    \item[Complex code retrieval] ColBERT or sparse retrieval may outperform pure dense for
          queries with specific identifiers, version numbers, or exact function names.
\end{description}
